\DocumentMetadata{
  pdfversion=2.0,
  pdfstandard=ua-2,
  lang=en-US,
  tagging=on,
  tagging-setup={math/setup=mathml-SE,tabsorder=structure}
}
\documentclass[11pt,letterpaper]{article}
\usepackage[margin=0.68in,includefoot]{geometry}
\usepackage{newcomputermodern}
\usepackage{amsmath,amscd,mathtools}
\usepackage{tikz,adjustbox}
\providecommand{\square}{\mdlgwhtsquare}
\usepackage[hidelinks]{hyperref}
\hypersetup{
  pdftitle={Algebra PhD Exam, January 2001},
  pdfauthor={University of Florida Department of Mathematics},
  pdfsubject={Graduate mathematics examination},
  pdfdisplaydoctitle=true,
  bookmarksopen=true
}
\setcounter{secnumdepth}{0}
\setlength{\parindent}{0pt}
\setlength{\parskip}{0.28em}
\setlength{\emergencystretch}{3em}
\setlength{\leftmargini}{2.15em}
\setlength{\leftmarginii}{2.65em}
\setlength{\leftmarginiii}{3em}
\pagestyle{empty}
\raggedbottom
\allowdisplaybreaks
\NewTaggingSocketPlug{block/list/label}{examproblem}{%
  \tagstructbegin{tag=Lbl}\tagstructend%
  \tagstructbegin{tag=\LBody}%
  \tagstructbegin{tag=\ProblemHeadingTag,title-o={\ProblemHeadingText}}%
  \tagmcbegin{tag=\ProblemHeadingTag,actualtext={\ProblemHeadingText}}%
  #2%
  \tagmcend%
  \tagstructend%
}
\newenvironment{examproblems}[2]{%
  \def\ProblemHeadingTag{#2}%
  \AssignTaggingSocketPlug{block/list/label}{examproblem}%
  \begin{enumerate}%
  \setcounter{enumi}{#1}%
  \renewcommand{\labelenumi}{\textbf{\theenumi.}}%
  \setlength{\topsep}{0.35em}%
  \setlength{\partopsep}{0pt}%
  \setlength{\itemsep}{0.48em}%
  \setlength{\parsep}{0.18em}%
}{\end{enumerate}}
\newenvironment{examsubparts}{%
  \AssignTaggingSocketPlug{block/list/label}{default}%
  \begin{enumerate}%
  \setlength{\topsep}{0.18em}%
  \setlength{\partopsep}{0pt}%
  \setlength{\itemsep}{0.16em}%
  \setlength{\parsep}{0.1em}%
}{\end{enumerate}}
\newenvironment{examinstructions}{%
  \par\begingroup\itshape
}{%
  \par\endgroup\vspace{0.18em}
}
\makeatletter
\renewcommand\subsection{\@startsection{subsection}{2}{0pt}%
  {-1.25ex plus -.25ex minus -.1ex}%
  {0.55ex plus .1ex}%
  {\normalfont\large\bfseries}}
\renewcommand{\@maketitle}{%
  \newpage\null\vskip -1em%
  {\footnotesize\bfseries
    \href{https://www.ufl.edu/}{University of Florida}, \href{https://math.ufl.edu/}{Department of Mathematics}\par}%
  \vskip -0.5em%
  \tagstructbegin{tag=H1,title={Algebra PhD Exam, January 2001}}%
  \tagpdfparaOff%
  \tagmcbegin{tag=H1}%
    {\LARGE\bfseries\href{https://gma.math.ufl.edu/exams/algebra-phd/}{Algebra PhD Exam}, January 2001\par}%
  \tagmcend%
  \tagpdfparaOn%
  \tagstructend%
  \vskip 0.25em%
  \tagmcbegin{artifact}%
    \hrule height 1pt%
  \tagmcend%
  \vskip 1em%
}
\makeatother
\title{Algebra PhD Exam, January 2001}
\author{}
\date{}
\begin{document}
\maketitle
\thispagestyle{empty}
\begin{examinstructions}
Time allowed: 240 minutes. Do seven of the following eleven problems. Please do not turn in more than seven problems. You must show your work. Answers with no work and/or no explanations will receive no credit. State clearly any theorem you use in your proofs. In the problems, \(\mathbb{Z}\), resp. \(\mathbb{Q}\), \(\mathbb{C}\), is the set of all integers, resp. of all rational numbers, of all complex numbers.
\end{examinstructions}
\begin{examproblems}{0}{H2}
\def\ProblemHeadingText{Problem 1.}
\item
Let~\(G\) be a finite simple (abelian or non-abelian) group of order~\(n\). Find the number of normal subgroups of \(G \times G\).
\def\ProblemHeadingText{Problem 2.}
\item
Show that there is no simple group of order \(112\). (Caution: You are not allowed to use Burnside’s \(p^a q^b\)-Theorem.)
\def\ProblemHeadingText{Problem 3.}
\item
Let~\(N\) be a minimal normal subgroup of a finite group~\(G\). Assume~\(N\) is solvable. Show that~\(N\) is an elementary abelian~\(p\)-group for some prime~\(p\).
\def\ProblemHeadingText{Problem 4.}
\item
Let \(A, B\) be two matrices over a field~\(F\). Suppose that \(\begin{pmatrix} A & 0 \\ 0 & A \end{pmatrix}\) is similar to \(\begin{pmatrix} B & 0 \\ 0 & B \end{pmatrix}\). Prove that~\(A\) is similar to~\(B\).
\def\ProblemHeadingText{Problem 5.}
\item
Determine the Galois group over~\(\mathbb{Q}\) of
\begin{examsubparts}
\item[\textnormal{a)}]
\((x^3-2)(x^2-3)\);
\item[\textnormal{b)}]
\((x^3-2)(x^2+3)\).
\end{examsubparts}
\def\ProblemHeadingText{Problem 6.}
\item
Is there any integer~\(d\) such that: for any separable extension~\(F\) of degree~\(5\) of a field~\(K\), if an extension~\(E\) of~\(K\) is a normal closure of \(F/K\) then \([E:K] \leq d\)? If yes, prove your bound. If not, justify why.
\def\ProblemHeadingText{Problem 7.}
\item
Let \(\xi \in \mathbb{C}\) be such that \(\xi^{2000}=3\). Show that \(-3\) is not a sum of squares in \(\mathbb{Q}(\xi)\).
\def\ProblemHeadingText{Problem 8.}
\item
Prove the Hilbert Basis Theorem: If~\(R\) is a commutative Noetherian ring with identity, then so is \(R[x_1,x_2,\ldots,x_n]\).
\def\ProblemHeadingText{Problem 9.}
\item
Let~\(A\) be a finite dimensional (not necessarily commutative)~\(\mathbb{C}\)-algebra with no zero divisors, but with identity. Show that \(\dim_{\mathbb{C}} A=1\).
\def\ProblemHeadingText{Problem 10.}
\item
This problem has two parts.
\begin{examsubparts}
\item[\textnormal{a)}]
If the ring~\(R\) is a principal ideal domain, prove that every nonzero prime ideal in~\(R\) is maximal.
\item[\textnormal{b)}]
Is the statement in a) still true if~\(R\) is not a principal ideal domain? (Hint: Consider the ring \(\mathbb{Z}[x]\).)
\end{examsubparts}
\def\ProblemHeadingText{Problem 11.}
\item
Suppose that~\(R\) is a ring with identity.
\begin{examsubparts}
\item[\textnormal{a)}]
Prove that each free left~\(R\)-module is projective.
\item[\textnormal{b)}]
Prove that a left~\(R\)-module~\(P\) is projective if and only if each short exact sequence of left~\(R\)-modules
\[
0 \longrightarrow A \longrightarrow B \longrightarrow P \longrightarrow 0
\]
 splits.
\end{examsubparts}
\end{examproblems}
\end{document}
